\documentclass{article}
\usepackage[T1]{fontenc}
\usepackage{lmodern}
\usepackage{graphicx}
\usepackage{amsmath,amssymb}
\usepackage[a4paper,margin=2cm]{geometry}
\usepackage[absolute,overlay]{textpos}
\setlength{\TPHorizModule}{1mm}
\setlength{\TPVertModule}{1mm}
\usepackage[indonesian]{babel}
\usepackage{hyperref}
\setlength{\emergencystretch}{2em}
% unit: Halaman 4

\begin{document}

% === Halaman 4 (python/native) ===

4

One-Time Pad (OTP)

• Satu-satunya algoritma kriptografi sempurna aman (perfect secrecy) sehingga 

tidak dapat dipecahkan adalah one-time pad (OTP).

• OTP ditemukan oleh Gilbert Vernam pada tahun 1917 dan kemudian secara 

matematis dibuktikan memiliki perfect secrecy oleh Claude Shannon tahun 

1949.

• OTP mengatasi kelemahan pada Vigenere Cipher. Vigenere Cipher mengulang 

penggunaan kunci secara periodik → mudah ditemukan dengan metode 

Kasiski.

• Pada OTP, panjang kunci = panjang plainteks, seluruh karakternya acak sejati

Plainteks: otpadalahcipheryangtidakbisadipecahkan

Kunci:       trjkdndkdwerylgrgdkopcegyhbdwjbtrfhgvk

\end{document}
